\section{Conservation is an exact statement about multiplicity}

The conservation identity is the least assumption-sensitive result in
this inventory \cite{yieldSource}. It is a finite-fibre count.

\begin{proposition}[Finite-pattern conservation]
Let $M$ and $G$ be finite sets and $s:M\to G$ a function. Put
$c(g)=|s^{-1}(g)|$. Then
\[
 \sum_{g\in G}c(g)=|M|.
\]
If $G$ is nonempty and the target is uniform, its mean multiplicity is
$|M|/|G|$.
\end{proposition}
\begin{proof}
Each element of $M$ belongs to exactly one fibre of $s$. The fibres
partition $M$, so their cardinalities sum to $|M|$. Division by the
number of targets gives the stated uniform mean.
\end{proof}

For a finite abelian group $G$, a set $D$ of $B$ distinct elements and
$m\geq1$, choose $M$ according to the declared counting convention:
\begin{align*}
 \text{ordered tuples with repetition:}\quad &|M|=B^m,\\
 \text{unordered multisets with repetition:}\quad
 &|M|=\binom{B+m-1}{m},\\
 \text{one element from each typed set }D_i:\quad
 &|M|=\prod_{i=1}^m |D_i|.
\end{align*}
The sum is well defined for multisets because the group is abelian.
Comparing two conventions without changing the denominator is not a
violation of conservation. Signed decompositions similarly require
specifying whether signs are separate labels or whether $D\cup(-D)$ is
a deduplicated set.

Coverage $p=|\{g:c(g)>0\}|/|G|$ is not the mean $\mu$.
Exactly,
\[
 p=\mu-\frac1{|G|}\sum_g(c(g)-1)_+\leq\min(1,\mu).
\]
When $\mu\leq1$, equality $p=\mu$ requires every count to be zero or one;
when $\mu>1$, equality with the bound $1$ requires coverage of every
target and permits repeated decompositions. Conservation therefore
leaves coverage, concentration and useful relation rank free. A claimed
universal numerical bound on a structured base's gain over a random base
would additionally require a model or bound for that random comparator;
it does not follow from this double count alone.

\paragraph{What is not conserved.}
The identity is not about the cost of finding a decomposition, the rank
of resulting rows, a nonuniform target law, or the mean over a selected
proper subset of $G$. A learned filter can redistribute work across
targets without contradicting it. A target-dependent base is another
experiment and must charge its selection. These are genuine open axes,
not exceptions to the stated identity.

The corpus already formalizes the general finite-pattern core in
\art{FactorBaseDesigns.sum_count}. It uses a list domain and a
nonduplicated target list covering its image, so domain repetitions
count as labelled multiplicities. The ordinary tuple or multiset
cardinality formula is a separate instantiation. Section~\ref{sec:formal}
records that boundary precisely.
